% concatenated from main.tex
\documentclass[12pt]{article}
\usepackage[utf8]{inputenc}
\usepackage[margin=0.8in]{geometry}
%\usepackage{tagstoc} %%% tag equations!

\usepackage{tocloft}


\usepackage{booktabs}\usepackage[svgnames]{xcolor}
\usepackage[bookmarksnumbered=true]{hyperref} 
\hypersetup{
     colorlinks = true,
     linkcolor = blue,
     anchorcolor = blue,
     citecolor = teal,
     filecolor = blue,
     urlcolor = blue
     }

   \usepackage{amsmath}
   \usepackage{amssymb}
   \newcommand\hmmax{0} % default 3
   % \newcommand\bmmax{0} % default 4
   \usepackage{bm}
\usepackage{soul}
%\usepackage{lmodern}
\usepackage{algorithm}
\usepackage{algpseudocode}
\usepackage{amsmath}
\usepackage{amssymb}
\usepackage{amsthm}
\usepackage{array}
\usepackage{wasysym}
\usepackage{mathtools}
\usepackage{graphicx}
\usepackage{hyperref}
\usepackage{booktabs}
\usepackage{bbold}
\usepackage{mathrsfs}
\usepackage{tikz-cd}
\usepackage{color}
\usepackage{setspace}
\usepackage{comment}
\usepackage{multirow}
\usepackage{bm}
\usepackage{thmtools}
\usepackage{thm-restate}
\usepackage{cleveref}
\usepackage{enumitem}
\usepackage{csquotes}
\usepackage{float}
\usepackage{authblk}
\usepackage{anyfontsize}


\usepackage{graphicx}
\usepackage{pgf,tikz,pgfplots}
\pgfplotsset{compat=1.17}
\usetikzlibrary{arrows}
\usetikzlibrary{angles,quotes}
\usetikzlibrary{calc}
\usepgfplotslibrary{fillbetween}
\usetikzlibrary{positioning}
\usetikzlibrary{calc,shapes.geometric}
\usetikzlibrary{decorations.pathreplacing}
\usetikzlibrary{cd}
\usetikzlibrary{positioning}
\usepackage{mathrsfs}
\usepackage{mathtools}
\usetikzlibrary{patterns}


\newcommand{\TT}{\mathbb{T}}
\newcommand{\NN}{\mathbb{N}}
\newcommand{\ZZ}{\mathbb{Z}}
\newcommand{\KK}{\mathbb{K}}
\newcommand{\QQ}{\mathbb{Q}}
\newcommand{\RR}{\mathbb{R}}
\newcommand{\CC}{\mathbb{C}}
\newcommand{\PP}{\mathbb{P}}
\newcommand{\BBA}{\mathbb{A}}


\newcommand{\FB}{\mathbf{B}}
\newcommand{\FD}{\mathbf{D}}
\newcommand{\fd}{\mathbf{d}}
\newcommand{\frv}{\mathfrak{v}}
\DeclareMathOperator{\rk}{rk}
\DeclareMathOperator{\Dr}{\rm Dr}
\DeclareMathOperator{\proj}{Proj}
\DeclareMathOperator{\Proj}{Proj}
\DeclareMathOperator{\Hom}{Hom}
\DeclareMathOperator{\Spec}{Spec}
\DeclareMathOperator{\Ima}{Im}
\DeclareMathOperator{\Trop}{Trop}

\newcommand{\MC}{\mathcal{C}}
\newcommand{\MI}{\mathcal{I}}
\newcommand{\MH}{\mathcal{H}}
\newcommand{\MX}{\mathcal{X}}
\newcommand{\MY}{\mathcal{Y}}
\newcommand{\ML}{\mathcal{L}}
\newcommand{\MM}{\mathcal{M}}
\newcommand{\MO}{\mathcal{O}}
\newcommand{\MP}{\mathcal{P}}
\newcommand{\MF}{\mathcal{F}}
\newcommand{\MS}{\mathcal{S}}
\newcommand{\mC}{\mathcal{C}}
\newcommand{\mY}{\mathcal{Y}}
\newcommand{\rank}{\textrm{Rank}}
%\newcommand{\Gr}{\textrm{Gr}}
%\newcommand{\Dr}{\textrm{Dr}}
\newcommand{\init}{\textrm{in}_{\prec}}
\newcommand{\Fl}{\mathscr{F}\!\ell}
\newcommand{\coker}{\mathrm{coker}}

\newcommand{\bigzero}{\mbox{\normalfont\Large\bfseries 0}}

\theoremstyle{plain}
\newtheorem{theorem}{Theorem}
\newtheorem{lemma}[theorem]{Lemma}
\newtheorem{corollary}[theorem]{Corollary}
\newtheorem{prop}[theorem]{Proposition}
\newtheorem{conjecture}[theorem]{Conjecture}
\newtheorem{question}{Question}
\newtheorem*{claim}{Claim}
\newtheorem{procedure}{Procedure}

\theoremstyle{definition}
\newtheorem{definition}[theorem]{Definition}
\newtheorem{remark}[theorem]{Remark}
\newtheorem{case}{Case}
\newtheorem{subcase}{Case}
\numberwithin{subcase}{case}
\newtheorem*{notation}{Notation}
\newtheorem{example}[theorem]{Example}




\declaretheorem[name=Conjecture]{conj}
\declaretheorem[name=Theorem]{thm}


%\theoremstyle{remark}

%\usepackage[backend=biber,
%style=alphabetic,
%sorting=nyt]{biblatex}
%\addbibresource{bibfile.bib}


\newcommand{\Bsig}{\mathcal{B}^\sigma}
\newcommand{\Btau}{\mathcal{B}^\tau}
\newcommand{\Bwo}{\mathcal{B}^{w_0}}
\DeclareMathOperator{\argmin}{argmin}
\DeclareMathOperator{\Gr}{Gr}
\DeclareMathOperator{\conv}{Conv}
\DeclareMathOperator{\Ker}{Ker}
\DeclareMathOperator{\vol}{Vol}
\DeclareMathOperator{\sgn}{sgn}
\DeclareMathOperator{\Sym}{Sym}
\usepackage{newtxmath}      % provides definitions for: \jmath, \amalg, \coprod
\hypersetup{
     colorlinks = true,
     linkcolor = blue,
     anchorcolor = blue,
     citecolor = teal,
     filecolor = blue,
     urlcolor = black
     }

\newcommand\restr[2]{{% we make the whole thing an ordinary symbol
  \left.\kern-\nulldelimiterspace % automatically resize the bar with \right
  #1 % the function
  \vphantom{\big|} % pretend it's a little taller at normal size
  \right|_{#2} % this is the delimiter
  }}

\usepackage{enumitem} % for '\newlist' and '\setlist' macros
\newlist{condenum}{enumerate}{1} % 'condenum': a new, enumerate-like list env.
\setlist[condenum]{label=\bfseries Condition \arabic*., 
                   ref=\arabic*, wide}
\newcommand\Myperm[2][^n]{\prescript{#1\mkern-2.5mu}{}P_{#2}}


\title{Generalized metric tree arrangements and Dressians
}
\author[1]{Ayush Kumar Tewari}
\affil[1]{Johann Radon Institute for Computational and Applied Mathematics (RICAM) \\
Austrian Academy of Sciences, Linz, Austria}
\affil[ ]{\texttt {tropical.tewari@gmail.com, ayushkumar.tewari@ricam.oeaw.ac.at}}
%\date{}




\begin{document}

\maketitle
\begin{abstract}
Metric trees and metric tree arrangements index cones in the polyhedral fan structure in the Dressian $\Dr(2,n)$ and $\Dr(3,n)$ respectively. We introduce the notion of generalized metric tree arrangements which parameterize points in $\Dr(k,n)$ and extend previously known results to $\Dr(k,n)$ along with providing explicit examples of these generalized metric tree arrangements. We study the adjacency of cones in the positive Dressian $\Dr_{>0}(3, n)$ and introduce generalized Whitehead moves which provide a condition for adjacency of maximal cones in $\Dr_{>0}(3,n)$ in terms of the associated metric tree arrangements. 
\end{abstract}

{\hypersetup{linkcolor=black}{\tableofcontents}}

%{\hypersetup{linkcolor=black}\setcounter{tocdepth}{1}\setlength\cftbeforesecskip{1.1pt}{\tableofcontents}}

\section{Introduction}

Tropical linear spaces are parameterized by the tropical prevariety Dressian $\Dr(k,n)$. The Dressian contains the tropicalization of the Grassmannian, which is the tropical Grassmannian $\Trop(\Gr(k,n))$. 

In \cite{speyer2004tropical} it is shown that $\Dr(2,n) = \Trop(\Gr(2,n))$ is a space which is equal to the \emph{moduli space of phylogenetic trees}. Additionally, in \cite{herrmann2009draw} the authors establish that the space of tropical planes $\Trop(\Gr(3,n))$ can be understood in terms of \emph{metric tree arrangement}. They compute explicitly $\Dr(3,6)$ and $\Dr(3,7)$ and describe all maximal cones, up to symmetry, and specify all the trivalent metric trees in the metric tree arrangement associated to each maximal cone. As tropical varieties enjoy a polyhedral fan structure, this perspective helps us understand the space $\Trop(\Gr(k,n))$ and the associated tropical prevariety i.e, Dressian $\Dr(k,n)$ as polyhedral fans. 

In this paper, we have two main objectives. Firstly, we try to understand the adjacency conditions for maximal cones in $\Dr_{>0}(3,n)$.  It is already known that metric trees indexing adjacent maximal cones in
$\Dr(2,n)$ are related by a \emph{Whitehead move}. Since. the polyhedral fan structure on the positive Dressian $\Dr_{>0}(3,n)$ is pure, which is not the case for the Dressian $\Dr(3,n)$ for $n \geq 7$, a well defined notion of adjacency of maximal cones exists for $\Dr_{>0}(3,n)$.
Building on the computations done on $\Dr(3,6)$ \cite{herrmann2009draw}, which is a pure polyhedral fan and hence also has adjacency well defined for maximal cones, we consider adjacent maximal cones in $\Dr(3,6)$ and introduce the notion of a \emph{generalized Whitehead move} in Definition \ref{def:gen_white_move}. We show that two adjacent maximal cones in $\Dr_{>0}(3,n)$ are related by a generalized Whitehead move in Theorem \ref{thm:Dr(3,n)_adjacency}. We formalize this result for $\Dr_{>0}(3,n)$ in the form of Lemma \ref{thm:Dr(3,n)_adj}.

%It is known that metric trees indexing adjacent maximal cones in $\Dr(2,n)$ are related by a \emph{Whitehead move}. 

We try to explore the connection between metric tree arrangements and $\Dr(k,n)$. Prior to this work, this is only known to exist for $\Dr(2,n)$ and $\Dr(3,n)$. We introduce the notion of \emph{generalized metric tree arrangements} and show that cones in $\Dr(k,n)$ index a collection of metric trees and also a generalized metric tree arrangement, and these two notions coincide in the case of $\Dr(3,n)$. This establishes the fact that our findings can be considered as a genuine extension of the results in \cite{herrmann2009draw} as illustrated by Corollary \ref{cor:final}. We believe this would open new avenues to study higher order Dressians, especially as we note that computing Dressians is computationally expensive \cite{herrmann2014dressians}, \cite{bendle2020parallel} and we stress on the fact that computations for $\Dr(4,8)$ are beyond the reach of current computational techniques and would require new techniques, although we do hope that with our work, for example the computations in Example \ref{eg:Dr(4,8)}; could provide some insights in developing these new techniques.

 
The structure of the paper is as follows. Section \ref{sec:Section_2} provides a brief introduction to the Dressian. In Section \ref{sec:Section_3} we delve into the question of adjacency for maximal cones in $\Dr_{>0}(3,n)$, staring with  adjacent maximal cones in $\Dr_{>0}(3,n)$ and proving the main technical result in Lemma \ref{thm:Dr(3,n)_adj}, based on which we state Theorem \ref{thm:Dr(3,n)_adjacency}, which is our final takeaway in terms of understanding adjacency in $\Dr_{>0}(3,n)$. Section \ref{sec:Section_4} introduces the definition of generalized metric tree arrangements and we show that these tree arrangements induce weights which reside in $\Dr(k,n)$ in Theorem \ref{prop:gen_tree_dressian}. We establish the equivalence of generalized metric tree arrangements and regular matroidal subdivisions of $\Delta(k,n)$, which connects our findings as the proper generalization of results concerning $\Dr(3,n)$ in \cite{herrmann2009draw}. Lastly, we have an appendix to provide descriptions of cones and their associated metric tree arrangements for $\Dr(3,6)$, based on which we formulate our results on $\Dr_{>0}(3,n)$. 


\section{Dressian}\label{sec:Section_2}

%\subsection{The structure of the Dressian Dr(2,n) }

We recall some basic definitions of Grassmannian, tropical Grassmannian and Dressian \cite{maclagan2021}.
The Grassmannian $\Gr(k,n)$ is a smooth projective variety defined by the Pl\"ucker ideal $I_{k,n}$,

\[ I _{k,n} = \langle  \mathcal{P}_{I,J} : I,J \subseteq [n], |I| = k-1, |J| = k+1 \rangle  \]

whose generators are the quadratic Pl\"ucker relations,

\[ \mathcal{P}_{I,J} = \sum_{j \in J} \text{sgn}(j;I,J) \cdot p_{I \cup j} \cdot p_{J \setminus j }  \]

The Grassmannian parameterizes $k$ dimensional subspaces of the vector space $\mathbb{K}^{n}$, where $\mathbb{K}$ is the underlying field. The \emph{tropical Grassmannian} $\Trop(\Gr(k,n))$ is the tropical variety defined by the Pl\"ucker ideal $I_{k,n}$. Among the generators of the ideal $I_{k,n}$ are the \emph{three term Pl\"ucker} relations which are of the form

\[ p_{Apq}p_{Ars} - p_{Apr}p_{Aqs} + p_{Aps}p_{Aqr}\]


where $A \in \binom{[n]}{k-2}$ and $p,q,r,s \in [n] \setminus A$ are pairwise distinct. The three-term Pl\"ucker relations generate the image of the ideal $I_{k,n}$ in the Laurent polynomial ring $\mathbb{K}[p_{A}^{\pm 1}]$, where  $A \in $ $\binom{[n]}{k-2}$. The \emph{Dressian} $\Dr(k,n)$ is the tropical prevariety defined by all the three-term Pl\"ucker relations. The Grassmannian and Dressian have a pure rational polyhedral fan structure in $\mathbb{R}^{{\binom{n}{k}} - 1 } \cong  \mathbb{R}^{\binom{n}{k}} / \mathbb{R}1$.

\begin{remark}
It is mentioned in \cite[Remark 12]{olarte2019local} that for the Dressian, considering all three term relations is equivalent to considering all quadratic relations and these two notions appear as strong and weak matroids over tropical hyperfields defined in \cite{baker2019matroids}, which are both equivalent.
\end{remark}

The Pl\"ucker relations define a natural Pl\"ucker fan structure on the Dressian $\Dr(k, n)$ where two weight vectors $\mu$ and $\mu'$ are in the same cone if they specify the same initial form for each relation. The Dressian $\Dr(k,n)$ is also endowed with a secondary fan structure as follows \cite{speyer2008linear_spaces},

\begin{theorem}
A vector $w \in \mathbb{R}^{ \binom{n}{k} }$ lies in the Dressian
$\Dr(k, n)$ if and only if it induces a matroidal subdivision of the hypersimplex $\Delta(k, n)$. 
\end{theorem}

It was shown in \cite{olarte2019local} that these two fan structures coincide.

\begin{theorem}\label{thm:coincide}
The Plücker fan structure on the Dressian $\Dr(k, n)$ as a fan in $\mathbb{R}^{{\binom{n}{k}} - 1 }$ coincides with the secondary fan structure.
\end{theorem}

This result is true in much more generality but Theorem \ref{thm:coincide} suffices for the purpose for our discussion.

We firstly discuss the case of $\Dr(2,n)$. $\Trop(\Gr(2,n))$ enjoys a strong connection with the space of phylogenetic trees which we recall here.

A \emph{phylogenetic} tree $T$ is a tree with $n$ labeled leaves and no vertices of degree two. A phylogenetic tree $T$ is \textit{trivalent} if for each vertex $v \in V(T)$ either $\deg(v) = 3$ or $\deg(v) = 1$. To each such phylogenetic tree we can attach metric data in the following way. We assign a length $l_{e} \in \mathbb{R}$ to each edge of $T$. Between any two leaves $i$ and $j$ of $T$, there is a unique path in $T$, and the distance between the leaves can be obtained as sum of $l_{e}$ along this path. This gives us the \emph{tree distance} vector $\delta = (\delta_{ij}) \in \mathbb{R}^{\binom{n}{2}}$. When all $l_{e}$ are nonnegative, the tree distance gives a finite metric. Finite metric spaces that arise from such metric tree are called \emph{tree metric}. The set of all tree distances is known as the space of phylogenetic trees on $n$ leaves and is denoted as $\mathcal{T}_{n}$ \cite{maclagan2021}.

\begin{theorem}[Theorem 4.3.5 \cite{maclagan2021}]\label{thm:rank2}
 $\Trop(\Gr(2,n)) = \Dr(2,n) = \mathcal{T}_{n}$
\end{theorem}

%In this section we recall the combinatorial structure of the maximal cells of the Dressian in terms of $3$-valent labelled trees \cite{speyer2004tropical}. For the Grassmannian $\Gr(2,n)$ we have that its tropicalisation $\Trop(\Gr(2,n))$ and the Dressian $\Dr(2,n)$ coincide. Moreover, the Dressian $\Dr(2,n)$ can be identified with the space of \textit{metric trees} with $n$ leaves. We proceed by recalling some basic notions that describe the connection between metric trees and the Dressian.







%The edges of a tree are said to be \textit{external} if they are incident to a leaf or \textit{internal} otherwise. A \textit{valid edge weight} on a on a tree $T$ is a function $\lambda : E(T) \rightarrow \RR$ such that $\lambda(e) > 0$ if $e \in E$ is an internal edge. A \textit{metric tree} is a pair $(T, \delta)$ where $T$ is a tree with $n$ leaves and $\delta : [n] \times [n] \rightarrow \RR$ is a function induced by some valid edge weight $\lambda$. Explicitly, if $i, j \in [n]$ are leaves of $T$ then $\delta(i,j)$ is the sum of all $\lambda(e)$ where $e$ lies in the path from $i$ to $j$.

%Each metric tree $(T, \delta)$ gives rise to a point in the Dressian as follows. Let $w \in \RR^{\binom{n}{2}}$ be the point whose $ij$'th coordinate is given by $w_{ij} = \delta(i,j)$. We have that $w$ is a point in the Dressian with respect to the maximum convention.

The space of metric trees naturally forms a polyhedral fan. The maximal cones of this fan correspond to trivalent trees. Two labeled trivalent trees $T_1$ and $T_2$ correspond to adjacent maximal cones if and only if they differ by a \textit{Whitehead move}. In Figure~\ref{fig:whitehead_move}, a Whitehead move is given by choosing an internal edge and modifying the tree as shown in the diagram. Note that the Whitehead move involves a choice since there are two possible ways to modify the trivalent tree.

\begin{figure}
    \centering
    \includegraphics[scale = 0.4]{Whitehead_move.drawio.pdf}
    \caption{Whitehead moves between three trivalent trees. The central figure shows the tree corresponding to the cone which is the intersection of the surrounding cones.}
    \label{fig:whitehead_move}
\end{figure}

We recall definitions of hypersimplex, splits, and associated results \cite{herrmann2008splitting}. The hypersimplex $\Delta(k,n)$ is the convex polytope obtained by intersecting the unit cube $[0,1]^{n}$ with the hyperplane $\sum x_{i} = d$. Equivalently, it can be defined as a $0-1$ polytope generated by the convex hull of all points with $k$ ones and $n-k$ zeros. The facet of $\Delta(k,n)$ defined by the hyperplane $x_{i} = 0$  is referred as the \emph{i-th deletion facet} and the facet defined by the hyperplane $x_{i} = 1$ is referred as the \emph{i-th contraction facet} \cite{herrmann2009draw}.

An important observation concerning deletion and contraction facets of $\Delta(k,n)$ is the following \cite{herrmann2009draw},

\begin{lemma}
Each deletion facet of $\Delta(k,n)$ is isomorphic to $\Delta(k,n-1)$ and each contraction facet of $\Delta(k,n)$ is isomorphic to $\Delta(d-1,n-1)$. 
\end{lemma}



\begin{definition}
A split $S$ of a polytope $P$ is the subdivision with exactly two maximal cells.
\end{definition}

The hyperplane defining the split $S$ is called the \emph{split hyperplane}.  A split is a regular subdivision \cite[Lemma 3.5]{herrmann2008splitting}. Two splits are said to be \emph{compatible} if the corresponding split hyperplanes do not meet in the interior of the polytope. 

Let $\Delta(k,n)$ be the hypersimplex. The $(A, B; \mu)$-hyperplane \cite{herrmann2008splitting} is defined as follows.
Let $A,B \subset [n]$ be non-empty sets such that $A \cup B = [n]$ and $A \cap B = \emptyset$. Let $\mu \in \mathbb{N}$. Then the linear equation
\[ \mu \sum_{i \in A} x_{i}  = (k - \mu) \sum_{i \in B} x_{i}\]
defines the $(A, B; \mu)$-hyperplane.

For any trivalent tree $T$ on $n$ leaves we consider an interior edge $e$. We collect the labels on either side of this edge $e$ in the sets $A$ and $B$. These define a $(A,B, \mu=1)$ split-hyperplane. 


The tree associated with a maximal cone gives a complete characterization of the matroids appearing in the subdivision of the hypersimplex.

\begin{theorem}
Let $T$ be a trivalent phylogenetic tree on $n$ leaves indexing a maximal cone of the Dressian $\Dr(2,n)$. The cells of the corresponding subdivision of the hypersimplex $\Delta(2,n)$ are indexed by the internal vertices of $T$. For each internal vertex $v \in V(T)$, the corresponding cell contains the basis $ij$ if and only if the path from $i$ to $j$ in $T$ passes through $v$.
\end{theorem}

\begin{proof}
From \cite{herrmann2008splitting} we know that a trivalent phylogenetic tree on $n$ leaves defines a compatible system of splits for $\Delta(2,n)$. We show that the matroidal subdivision that is obtained from the common refinements of this system of splits coincides with the matroidal subdivision obtained from the indexing of the trivalent tree $T$.

We consider the labelling $a_{1} \hdots a_{n}$ on $T$. Every internal edge of $T$ provides a $(A,B)$ split of the form 

\begin{align*}
(A,B)_{e_{1}} = (\{a_{1},a_{2}\}, \{a_{3}, \hdots \}) \\
\hdots                          \hdots \\
\hdots                          \hdots \\
(A,B)_{e_{n-3}} = (\{a_{1}, \hdots a_{n-2}\}, \{a_{n-1}, a_{n}\}) \\
\end{align*}

Each split provides two matroid polytopes of the following form 

\[ (S_{1})_{e_{1}} = \{ a_{1}a_{2}, a_{1}a_{3}, \hdots , a_{1}a_{n}, a_{2}a_{3}, \hdots a_{2}a_{n}\}
\]
\[ (S_{1})_{e_{2}} = \{ a_{1}a_{3}, \hdots  a_{1}a_{n}, a_{2}a_{3}, \hdots a_{2}a_{n}, \hdots a_{n}a_{n-1}, \hdots a_{3}a_{n} \}
\]

The common refinements of all these n-3 splits provides us the following matroidal subdivision $\Sigma$ with $n-2$ maximal cells,

\begin{align*}
    M_{1} = \{ a_{1}a_{2}, \hdots , a_{2}a_{n} \}
    \hdots                    \\
    \hdots \\
    M_{n-2} = \{ a_{n-1}a_{1}, \hdots , a_{n}a_{2} \}
\end{align*}

As we see these maximal cells are indexed by the $n-2$ internal vertices of $T$. Hence, we see that $\Sigma$ coincides with the matroidal subdivision obtained by the indexing of the tree $T$.
%\begin{proof}
    %No idea. See Figure~\ref{fig:tree_to_matroid}.
%\color{red}
%One possible approach is to show that the trivalent labeled tree $T$ is dual to the %corresponding matroidal subdivision, as a a polytopal subdivision of $\Delta(2,n)$ is a matroid
%subdivision if and only if it is dual to a tree \cite[Lemma 4.2]{herrmann2009draw}. 
%\color{black}
    %\begin{figure}[h]
        %\centering
        %\includegraphics[scale = 0.8]{tree_to_matroid_subdivision.pdf}
        %\caption{Two trees for $\Gr(2,6)$ that differ by a Whitehead move together with their associated matroid subdivisions. The paths associated to the basis $14$ are highlighted in each tree.}
        %\label{fig:tree_to_matroid}
    %\end{figure}
\end{proof}

%\medskip


\medskip

We illustrate this with the following example for $\Dr(2,5)$.   



\begin{example}[$\Dr(2,5)$]
We take the tree $T$ on 5 leaves as shown in Figure \ref{fig:tree_to_matroidal_subdiv}. The two internal edges $e_{1}$ and $e_{2}$ define the following splits,
\[(A,B)_{e_{1}} = \{\{1,2\},\{3,4,5\}\} \]
\[(A,B)_{e_{2}} = \{\{1,2,3\},\{4,5\}\} \]
Since the points of the matroid polytopes reside on either side of the hyperplanes defined by the corresponding split, we see that the matroid polytopes defined by the split $(A,B)_{e_{1}}$ are
\[ S_{1} = \{12,13,14,15,23,24,25\} \]
\[ S_{2} = \{13,14,15,23,24,25,34,35\} \]
Similarly, for the split $(A,B)_{e_{2}}$
\[ S'_{1} = \{14,15,24,25,34,35,45\} \]
\[ S'_{2} = \{14,15,24,25,34,35,13,23\} \]
From \cite{herrmann2008splitting} we know that such a system of splits obtained from a tree are compatible, and the common refinement of these two splits provides us the following regular matroidal subdivision
\[ M_{v_{1}} = \{12,13,14,15,23,24,25\} \]
\[ M_{v_{2}} = \{13,14,15,23,24,25,34,35\} \]
\[ M_{v_{3}} = \{14,15,24,25,34,35,45\} \]
which is a regular matroidal subdivision of $\Delta(2,5)$ whose maximal cells are indexed by the internal vertices of the tree $T$ nad lists all paths passing through the respective internal vertex.
 \begin{figure}[h]
        \centering
        \includegraphics[scale = 0.5]{tree_example.drawio.pdf}
        \caption{A trivalent tree for Dr(2,5)}
        \label{fig:tree_to_matroidal_subdiv}
    \end{figure}
\end{example}

\begin{definition}

Let $n \geq 4$ and we consider a $n$-tuple of metric trees $T = (T_{1}, \hdots , T_{n})$, each with an associated metric $\delta_{i}$. We call the $n$- tuple $T$ of metric trees a \emph{metric tree arrangement} if, 

\[ \delta_{i}(j,k) = \delta_{j}(i,k) =\delta_{k}(i,j) \]

for all $i,j,k \in [n]$ pairwise distinct \cite{herrmann2009draw}.

Moreover, considering trees $T_{i}$ without metrics, but
with leaves still labeled by $[n] \setminus {i}$, we say that $T$ is an \emph{abstract tree arrangement} if,

\begin{itemize}
    \item $n =4$,
    \item $n = 5$, and $T$ is the set of quartets of a tree with five leaves,
    \item $n \geq 6$, and $(T_{1} \setminus i, \hdots , T_{i-1} \setminus i, T_{i+1} \setminus i, \hdots , T_{n} \setminus i)$ is an arrangement of $n - 1$ trees for each $i \in [n]$ \cite{herrmann2009draw}.
\end{itemize}

\end{definition}

Metric tree arrangements index cones of $\Dr(3,n)$, and this is elaborated by the following results in \cite{herrmann2009draw}.

\begin{lemma}[Lemma 4.2 in \cite{herrmann2009draw}]
Each matroid subdivision $\Sigma$ of $\Delta(3, n)$ defines an abstract arrangement $T(\Sigma)$ of $n$ trees. Moreover, if $\Sigma$ is regular then $T(\Sigma)$ supports a metric tree arrangement.
\end{lemma}

Two metric tree arrangements are \emph{equivalent} if they induce the same abstract tree arrangement.

\begin{theorem}[Theorem 4.4 \cite{herrmann2009draw} ]
The equivalence classes of arrangements of $n$ metric trees are in bijection with the regular matroid subdivisions of the hypersimplex $\Delta(3,n)$. 
\end{theorem}
 

Before we progress to the next section, we also recall some definitons concerning polyhedral fans \cite{arkani2021positive}. 

\begin{definition}
A polyhedral fan $\mathcal{F} = \{C\}$ in $\mathbb{R}^{n}$ is a finite collection of polyhedral cones $C \subset \mathbb{R}^{n}$ such that
\begin{itemize}
    \item if $C, C' \in \mathcal{F}$, then $C \cap C' \in \mathcal{F}$
    \item for $C \in \mathcal{F}$, every face $C'\subset C$ is in $\mathcal{F}$
\end{itemize}
\end{definition}

The fan $\mathcal{F}$ is called \emph{complete} if the union of all cones in it is $\mathbb{R}^{n}$ and $\mathcal{F}$ is called \emph{pure} if all maximal cones have the same dimension. Two maximal cones $C$ and $C'$ of a pure polyhedral fan $\mathcal{F}$ of dimension $d$ are said to be \emph{adjacent} if they intersect in a common face $F$ of dimension $d-1$. An important thing to note here is that the notion of adjacency for cones is not well defined for polyhedral fans that are not pure.

We now recall two other subspaces of the Grassmannian $\Gr(k,n)$, namely the \emph{totally nonnegative Grassmannian} $\Gr_{\geq 0}(k,n)$ and the \emph{positive Grassmannian} $\Gr_{> 0}(k,n)$ both of which were introduced by Lustig \cite{lusztig1994total} and Postnikov \cite{postnikov2006total}. A point $V \in \Gr_{\geq 0}(k,n)$ if all its Plücker coordinates are nonnegative. Similarly, a point $V \in \Gr_{> 0}(k,n)$ if all its Plücker coordinates are positive. The underlying matroid for a point residing in the totally nonnegative Grassmannian is a \emph{positroid} and the totally nonnegative Grassmannian admits a stratification in terms of positroid cells. We know that points residing in the Dressian $\Dr(k,n)$ satisfy the \emph{tropical Plücker relations} \cite{speyer2008linear_spaces}. Similar, to this the authors in \cite{arkani2021positive} introduce the \emph{positive tropical Plúcker relations}. A point $p = \{p_{I} | I \in \binom{n}{k} \}$ where $p_{I} \in \mathbb{R} \cup \{\infty\}$ and not all $p_{I}$ are infinity, is said to satisfy a \emph{positive tropical Plücker} relation if for every set $S$ of size $k-2$ and $a < b < c < d$,

\[ p_{Sac} + p_{Sbd} = \text{min}(p_{Sab} + p_{Scd}, p_{Sad} + p_{Sbc}) \]

The support of the point $p$ is defined as $\text{Supp}(p) = \{ I | p_{I} < \infty \}$. For a positroid $M$, let $\Dr(M)_{ > 0}$ denote the set of positive tropical Plücker vectors and we refer to it as the \emph{positive Dressian} with support $M$. $\Dr(k,n)_{ > 0}$ is the positive Dressian when the underlying matroid is the uniform matroid $U(k,n)$, which is also a positroid.



By the structure theorem of tropical varieties \cite{maclagan2021introduction}, we know that $\Trop(\Gr(k,n))$ has the structure of a pure polyhedral fan for all values of $k$ and $n$. Therefore, as a consequence of Theorem \ref{thm:rank2} Dressian $\Dr(2,n)$ is pure for all $n$. However, it is also known that the Dressian $\Dr(3,n)$ is not a pure polyhedral fan for $n \geq 7$ \cite{herrmann2009draw}. But if we restrict ourselves to the positive Dressian $\Dr_{> 0}(3,n)$, we have an interesting observation in the form of the following result,

\begin{theorem}[Proposition 10.2 \cite{arkani2021positive}]\label{thm:dress_pure}
The positive Dressian $\Dr_{> 0}(3,n)$ is a pure polyhedral fan.
\end{theorem}

Additionally, it also known that in this case the different fan structures coincide,

\begin{theorem}[Theorem 10.3 \cite{arkani2021positive}]\label{thm:pos_dress_coincide}
The secondary fan structure and the Plücker fan structure coincide for the positive $\Dr_{> 0}(3,n)$.    
\end{theorem}


\begin{remark}
Both Theorem \ref{thm:dress_pure} and Theorem \ref{thm:pos_dress_coincide} have been proven in much more generality in \cite{arkani2021positive} for the positive Dressian $\Dr_{> 0}(M)$ of an arbitrary positroid $M$.    
\end{remark}

\section{Adjacency of maximal cones in positive Dressian \texorpdfstring{$\Dr_{> 0}(3,n)$}{Dr(3,n)}}\label{sec:Section_3}



%\subsection{Adjacency in cones in Dr(3,n)}

We describe some notation for our results. Let $\Sigma^{\alpha}$ be a regular matroidal subdivision indexing a cone $C$ in $\Dr(3,n)$. We refer $T^{\alpha}$ as the metric tree arrangement for the cone $C$. The matroidal subdivision induced on the i-th contraction facet $\{x_{i} = 1\}$ by $\Sigma^{\alpha}$ is referred as $\Sigma^{\alpha}_{i}$.

\begin{lemma}\label{thm:Dr(3,n)_adj}
 Let $\Sigma^{\alpha}$ and $\Sigma^{\beta}$ be two regular positroidal subdivisions of $\Delta(3,n)$ such that $\Sigma^{\alpha}$ and $\Sigma^{\beta}$ lie in adjacent maximal cones of $\Dr_{> 0}(3,n)$. Let $T^{\alpha}$ and $T^{\beta}$ be the metric tree arrangements corresponding to $\Sigma^{\alpha}$ and $\Sigma^{\beta}$. Then, $\forall i \in [n]$ either, 
 \begin{enumerate}
     \item $\Sigma^{\alpha}_{i} = \Sigma^{\beta}_{i}$, or
     \item $\Sigma^{\alpha}_{i}$ and $\Sigma^{\beta}_{i}$ lie in adjacent maximal cones of $\Dr(2,n-1)$
 \end{enumerate}
\end{lemma}

\begin{proof}
By \cite{herrmann2009draw}[Lemma 4.3] we know that each regular matroidal subdivision $\Sigma$ of $\Delta(3,n)$ induces a regular matroidal subdivision $\Sigma_{i}$ on the i-th contraction facet isomorphic to $\Delta(2,n-1)$. Let $T^{\alpha}_{i}$ be the tree in $T^{\alpha}$ which is dual to the subdivision $\Sigma^{\alpha}_{i}$. 

We proceed by contradiction and assume that there exists a set $J \subsetneq [n]$, such that 

\[ \Sigma^{\alpha}_{j} = \Sigma^{\beta}_{j}, \quad \forall j \in J \]

but for all $j' \in [n] \setminus J$,  $\Sigma^{\alpha}_{j'}$ and $\Sigma^{\beta}_{j'}$ are not adjacent to each other, i.e., there exists at least one cone $\Sigma^{\gamma}_{j"}$ (for some $\Sigma^{\gamma} \in \Dr_{> 0}(3,n)$) which lies in between  $\Sigma^{\alpha}_{j'}$ and $\Sigma^{\beta}_{j'}$ in $\Dr(2,n-1)$. Since, adjacency in $\Dr(2,n-1)$ is determined by a Whitehead move, hence $j" = j'$, as all trees related by a Whitehead move have the same set of leaves. 

We now define a metric tree arrangement $T^{\gamma}$ such that 

\[ T^{\alpha}_{j} = T^{\beta}_{j} = T^{\gamma}_{j}, \quad \forall j \in J \]

and for all $j' \in [n] \setminus J$, the tree $T^{\alpha}_{j'}$ is related by a Whitehead move to $T^{\gamma}_{j'}$ and $T^{\gamma}_{j'}$ is related to $T^{\beta}_{j'}$ by a Whitehead move. By construction, $T^{\gamma}$ is not equivalent to the abstract tree arrangements corresponding to $T^{\alpha}$ and $T^{\beta}$. Hence, by \cite{herrmann2009draw}[Theorem 4.5], $T^{\gamma}$ corresponds to the regular matroidal subdivision $\Sigma^{\gamma}$. We now show that $\Sigma^{\gamma}$ is in between $\Sigma^{\alpha}$ and $\Sigma^{\beta}$ to arrive at our desired contradiction.

In \cite{herrmann2009draw}, it is shown that for any metric tree arrangement $T$ in $\Dr(3,n)$, we can assign the following map 

\begin{equation*}
    \pi(i,j,k) =
    \begin{cases}
      \delta_{i}(j,k) = \delta_{j}(i,k) = \delta_{k}(i,j), & \text{if}\ \text{i,j and k pairwise distinct} \\
      \infty & \text{otherwise}
    \end{cases}
  \end{equation*}

to increasing tuples $i <j < k$ to obtain the finite tropical Pl\"lucker vector which lies in the Pl\"ucker fan of $\Dr(3,n)$. Hence, for the tree arrangements $T^{\alpha}, T^{\beta}$ and $T^{\gamma}$ we obtain the vectors $\pi^{\alpha}, \pi^{\beta}$ and $\pi^{\gamma}$. For $\pi^{\alpha} \in \Dr(3,n)$, this implies,

\[\text{min}\{\pi^{\alpha}_{hij} + \pi^{\alpha}_{hkl}, \pi^{\alpha}_{hik} + \pi^{\alpha}_{hjl}, \pi^{\alpha}_{hil} + \pi^{\alpha}_{hjk}  \} \]

is attained at least twice for all pairwise distinct $h,i,j,k,l \in [n]$. But this Pl\"ucker relation can be seen as the Pl\"ucker relation for the tree $T_{h} \in T_{\alpha}$

\[ \text{min} \{\pi^{\alpha}_{ij} + \pi^{\alpha}_{kl}, \pi^{\alpha}_{ik} + \pi^{\alpha}_{jl}, \pi^{\alpha}_{il} + \pi^{\alpha}_{jk}  \} \]

which gives a point in $\Dr(2,n-1)$. So for the $n$ trees in $T^{\alpha}$ we have $\binom{n}{1}$ choices of $h$, and for each such choice we have $\binom{n}{4}$ Pl\"ucker relations in $\Dr(2,n-1)$.  We know that adjacency in $\Dr(2,n-1)$ is determined by Whitehead moves, so for all $j' \in [n] \setminus J$, the initial forms corresponding to Pl\"ucker relations for tree $T_{j'}$ lie in adjacent maximal cones of $\Dr(2,n-1)$ for $\pi^{\alpha}$ and $\pi^{\gamma}$. By replacing $h$ in the indices we obtain the corresponding initial forms for Pl\"ucker relations which determine points in adjacent maximal cones in $\Dr_{> 0}(3,n)$. Hence, we conclude that $\pi^{\alpha}$ and $\pi^{\gamma}$ lie in adjacent maximal cones of the Pl\"ucker fan in $\Dr_{> 0}(3,n)$. Similarly, $\pi^{\beta}$ and $\pi^{\gamma}$ lie in adjacent maximal cones of the Pl\"ucker fan in $\Dr_{> 0}(3,n)$. Therefore, the maximal cone corresponding to $\pi^{\gamma}$ is in between $\pi^{\alpha}$ and $\pi^{\beta}$. By Theorem \ref{thm:pos_dress_coincide} we know that the Pl\"ucker fan structure and secondary fan structure on the positive Dressian $\Dr_{> 0}(3,n)$ coincide which implies that the cone for $\Sigma^{\gamma}$ is in between $\Sigma^{\alpha}$ and $\Sigma^{\beta}$, which gives us the desired contradiction.

\end{proof}

Hence, we now know the relation between tree arrangements of adjacent maximal cones in $\Dr_{>0}(3,n)$. With this, we now introduce the definition of \emph{generalized Whitehead moves}. 

\begin{definition}\label{def:gen_white_move}
Consider two metric tree arrangements $T^{\alpha}$ and $T^{\beta}$ containing $m$ trees. Let $T^{\alpha}_{i}$ and $T^{\beta}_{i}$ denote the $i$-th metric tree respectively.  We say that $T^{\alpha}$ and $T^{\beta}$ are related via generalized Whitehead move if, for a set $J \subset [m]$, 

\begin{itemize}
    \item $T^{\alpha}_{j} = T^{\beta}_{j} \quad \forall j \in J$
    \item $T^{\alpha}_{j'}$ and $T^{\beta}_{j'}$ is related by a Whitehead move $\forall j' \in [m] \setminus J$ 
\end{itemize}
\end{definition}

With this definition, we state the following result,

\begin{theorem}\label{thm:Dr(3,n)_adjacency}
Two adjacent maximal cones in $\Dr_{>0}(3,n)$ are related by a generalized Whitehead move.
\end{theorem}

\begin{proof}
The proof follows from Lemma \ref{thm:Dr(3,n)_adj} and Definition \ref{def:gen_white_move}.
\end{proof}


Theorem \ref{thm:Dr(3,n)_adjacency} agrees with our computational observations for $\Dr(3,6)$, which is known to be a pure polyhedral fan \cite{herrmann2009draw}. In this case, adjacency of cones is well defined not only for the positive Dressian $\Dr_{>0}(3,6)$ but for the whole Dressian $\Dr(3,6$), and hence we compute adjacent maximal cones in the Dressian $\Dr(3,6)$. This computation also illustrates that the roadblock for an extension of Theorem \ref{thm:Dr(3,n)_adjacency} to the case of the Dressian $\Dr(3,n)$ is the fact that it is a non-pure polyhedral fan for $n \geq 7$, because of which a adjacency notion for maximal cones is not well defined. In Appendix \ref{app:appendix}, we see that adjacent maximal cones in $\Dr(3,6)$ have trees that are either identical or related via a Whitehead move. This can be seen evidently in Figure \ref{fig: cone_in_bw_cone9andcone10}, where we look at the tree arrangement on the wall between the adjacent maximal cones Cone$_{9}$ and Cone$_{10}$ in $\Dr(3,6)$ which is spanned by the ray generators $\{42,53,54\}$ and we see that the trees are related by Whitehead moves, which shows up as trees with a degree four vertex, which can be understood as the collapse of one of the internal edges involved in the Whitehead move. 

\begin{figure}
    \centering
    \includegraphics[scale=0.8]{tree_on_codimension1_face.drawio.pdf}
    \caption{Tree arrangement corresponding to the wall in between adjacent  cones Cone$_{9}$ and Cone$_{10}$ in $\Dr(3,6)$}
    \label{fig: cone_in_bw_cone9andcone10}
\end{figure}

We now point the reader to some questions regarding the adjacency of maximal cones in Dressian $\Dr_{> 0}(3,n)$ that we would like to tackle in the future,

\begin{itemize}
    \item Can the characterization of adjacency of cones in terms of Whitehead moves between certain trees be extended to the positive Dressian $\Dr_{> 0}(k,n)$ for $k \geq 4$? In the next section, with the introduction of generalized metric tree arrangements for $\Dr(k,n)$, a setup to try to prove such a result is provided.
    
    \item Can we bound the number of trees related via Whitehead move between two adjacent maximal cones for $\Dr_{> 0}(3,n)$? In our example of $\Dr(3,6)$, this bound turns out to be three.
    
    \item Can we obtain indices of the trees which are related via Whitehead moves between adjacent maximal cones, from only the information of the cones?  
\end{itemize}





\section{Generalized metric tree arrangements for Dr(k,n)}\label{sec:Section_4}

We begin this section with a discussion on the structure of Dressian  
$\Dr(k,n)$ ($k \geq 4)$. Firstly, we state Lemma \ref{lem:mat_sub_k,n} which is a generalization of \cite[Lemma 4.3]{herrmann2009draw},

\begin{lemma}\label{lem:mat_sub_k,n}
Each matroid subdivision $\Sigma$ of $\Delta(k,n)$ defines an abstract tree arrangement $T(\Sigma)$ of $\Myperm[n]{k-2}$ trees, with each tree having $n-(k-2)$ leaves. Moreover, if $\Sigma$ is regular then $T(\Sigma)$ supports metric tree arrangement.
\end{lemma}

\begin{proof}
Each of the $n$ contraction facets of $\Delta(k,n)$ are  isomorphic to $\Delta(k-1,n-1)$. Recursively using this relation we get that $\Sigma$ induces matroidal subdivisions on $(k-2)! \cdot {\binom{n}{k-2}} = \Myperm[n]{k-2}$ copies of $\Delta(2,n-(k-2))$. But matroidal subdivisions of $\Delta(2,n-(k-2))$ are generated by the compatible system of splits, and such subdivisions are dual to trees. Hence, $\Sigma$ gives rise to a tree arrangement.

Let $\Sigma$ be regular with $\lambda$ as a weight function. We now invoke arguments, similar to what has been used in proof of Lemma 4.2 in \cite{herrmann2009draw}. By suitable rescaling, we can assume that $\lambda$ has values between $1$ and $2$. We restrict $\lambda$ to $\Delta(2,n-(k-2))$ to obtain the induced subdivisions. But a weight function on $\Delta(2,n-(k-2))$ which takes values between $1$ and $2$ is a metric. Since, the induced subdivisions on $\Delta(2,n-(k-2))$ are also regular matroidal subdivisions, they are dual to metric trees on $n-(k-2)$ leaves. The Split Decomposition theorem \cite[Theorem 2.11]{herrmann2008splitting} provides the lengths on all edges of these trees.
\end{proof}

\begin{example}[Metric tree arrangement for $\Delta(4,8)$]\label{eg:Dr(4,8)}
We now describe the metric tree arrangement supported by a regular matroidal subdivision of $\Delta(4,8)$. We obtain the following matroidal subdivision by a lattice path construction,


\begin{equation*}
\begin{split}
 M_{1} & = \{0,1,2,3,4,5,6,7,8,15,16,17,18,35,36,37,38\} \\
 M_{2} & = \{1,2,3,4,5,6,7,8,9,10,11,15,16,17,18,19,20,21,35,36,37,38,39,40,41\} \\
  M_{3} & = \{5,6,7,8,9,10,11,15,16,17,18,19,20,21,25,26,27,35,36,37,38,39,40,41,45,46,47\} \\
  M_{4} &= \{6,7,8,9,10,11,12,13,14,16,17,18,19,20,21,22,23,24,25,26,27,28,29 \\
 & 30,31,32,33,34,36,37,38,39,40,41,42,43,44,45,46,47,48,49,50 \\
  & 51,52,53,54,55,56,57,58,59,60,61,62,63,64,65,66,67,68,69\} \\
   M_{5} & = \{15,16,17,18,25,26,27,35,36,37,38,45,46,47,55,56,57\} 
\end{split}
\end{equation*}

Using {\tt polymake} \cite{polymake:2000}, we find that this subdivision is induced by the following weight vector 

\begin{equation*}
\begin{split}
w & =\{0,0,0,0,0,0,0,0,0,1,1,1,2,2,2,0,0,0,0,1,1,1,2,2,2,2,2,2, \\
 & 3,3,3,4,4,4,5,0,0,0,0,1,1,1,2,2,2,2,2,2,3,3,3, \\
 & 4,4,4,5,2,2,2,3,3,3,4,4,4,5,5,5,5,6,7\} 
\end{split}
\end{equation*}



The subdivision induced by this weight vector on the contraction facet isomorphic to $\Delta(3,7)$ is given by the hyperplane of the form $\{x_{i} = 1\}$, and for $x_{1} = 1$, the induced weight vector is 

\begin{equation*}
\begin{split}
 w^{\text{in}} & = \{0,0,0,0,0,0,0,0,0,1,1,1,2,2,2,0,0,0,0,1,
  \\
  & 1,1,2,2,2,2,2,2,3,3,3,4,4,4,5\} 
\end{split}
\end{equation*}


\begin{figure}[H]
    \centering
    \includegraphics[scale=0.31]{mat_4,8_example.drawio.pdf}
    \caption{The tree arrangement for subdivision induced on a contraction facet of $\Delta(4,8)$ together with the maximal cells of the corresponding matroid subdivision.}
    \label{fig: mat(4,8)}
\end{figure}



With respect to the labeling of leaves with $[1,2 \hdots 7]$ Figure \ref{fig: mat(4,8)} shows the corresponding tree arrangement to this induced matroidal subdivision, along with the corresponding maximal cells for the subdivision. The full metric tree arrangement for the matroidal subdivision of $\Delta(4,8)$ is obtained by considering all the other contraction facets and obtaining the corresponding tree arrangements on each copy of $\Delta(3,7)$. We also note that this subdivision is clearly coarser than a subdivision induced by a point residing in the maximal cone of $\Dr(4,8)$, hence we see that not all trees in the tree arrangement are trivalent.


\end{example}

We now describe an extension of the definition of metric tree arrangement \cite{herrmann2009draw}. 

\begin{definition}\label{def:gen_metric_tr_arr}
A \textit{generalized metric tree arrangement for $\Dr(k,n)$} is a set $T = \{T_I \colon I \subseteq [n], \ |I| = k-2 \}$ of metric trees such that the leaves of $T_I$ are labeled by the elements of $[n] \backslash I$ and $T$ satisfies the following compatibility condition: for any pair of metric trees $T_I$ and $T_{I'}$ in $T$; any pair of leaves $i, j$ of $T_I$; and any pair of leaves $i', j'$ of $T_{I'}$, if $I \cup \{i,j\} = I' \cup \{i', j'\}$ then we have

\begin{equation}\label{eq:comptability}
 \delta_I(i,j) = \delta_{I'}(i', j')
\end{equation}

where $\delta_I$ is the metric on $T_I$ and $\delta_{I'}$ is the metric on $T_{I'}$. We define $\TT(k,n)$ to be the collection of all generalized metric tree arrangements for $\Dr(k,n)$.
\end{definition}


\begin{remark}
It is evident from the definition that a metric tree arrangement is also a generalized metric tree arrangement and corresponds to the case when $k=3$.
\end{remark}


Each generalised metric tree $T$ induces a weight vector $w_T$ on the Pl\"ucker variables given by $w_T(P_K) = \delta_J(i,j)$ where: $K = J \cup \{i,j\}$; $\delta_J$ is the metric on the tree $T_J$; and $i,j$ are leaves of the tree $T_J$. Note that the definition of the metric tree arrangement above ensures that $w_T$ is a well-defined weight on the Pl\"ucker variables. We establish the relation between the weight vector $w_{T}$ and the Pl\"ucker fan structure of $\Dr(k,n)$.

\begin{theorem}\label{prop:gen_tree_dressian}
As sets, the collection of weights $\{w_T \colon T \in \TT(k,n) \}$ is equal to the Dressian $\Dr(k,n)$.
\end{theorem}

\begin{proof}
Let $T \in \TT(k,n)$ be a tree arrangement. Fix a $3$-term tropical Pl\"ucker relation for $\Dr(k,n)$. Note that this relation is of the form
\[
P = 
P_{J \cup ij} \odot P_{J \cup k\ell} \oplus
P_{J \cup ik} \odot P_{J \cup j\ell} \oplus
P_{J \cup i\ell} \odot P_{J \cup jk} 
\]
where $J \subseteq [n]$ has size $k - 2$ and $i,j,k,\ell \in [n] \backslash J$ are distinct. Let $S = \{I \subseteq [n] \colon |I| = k, \ J \subseteq I \}$ be the $k$-subsets containing $J$. By definition, the restriction of $w_T$ to the Pl\"ucker variables $P_I$ where $I \in S$ is induced by a single metric tree $T_J \in T$. By \cite{speyer2004tropical}, the space of weights induced by metric trees with $n-k+2$ leaves equals the Dressian $\Dr(2, n-k+2)$. Therefore the maximum in $P$ is attained by at least two terms. So we have shown that $\{w_T \colon T \in \TT(k,n) \} \subseteq \Dr(k,n)$.

For the other inclusion take any point $w \in \Dr(k,n)$. For each $(k-2)$-subset $J \subseteq [n]$, we construct a metric tree $T_J$ with leaves $[n]\backslash J$. Consider the restriction of $w$ to the Pl\"ucker variables $P_I$ where $J \subseteq I$. Note that these Pl\"ucker variables are naturally in bijection with the $2$-subsets of $[n]\backslash J$. So, $w$ gives rise to point in $\Dr(2, n-k+2)$ given by $\hat w(P_{ij}) = w(P_{J \cup ij})$. As above, the Dressian $\Dr(2, n-k+2)$ is equal to the space of weights induced by metric trees and so we define $T_J$ to be the metric tree associated to $\hat w$ with leaves labeled by $[n] \backslash J$. It remains to show that the trees satisfy the compatibility condition. However, this is immediate since each metric tree $T_J$ is constructed from the same point $w \in \Dr(k,n)$.
\end{proof}

Alternatively, the result in Theorem \ref{prop:gen_tree_dressian} can be understood as a generalization of the map $\pi$ defined in the proof of Theorem 4.4 in \cite{herrmann2009draw}. We describe this perspective here. Let $\pi$ be the map from $[n]^{k} \rightarrow \mathbb{R} \cup \{\infty\} $

\[ \pi(i_{1}, i_{2}, \hdots , i_{k}) = \begin{cases}
  \delta_{I_{1}}(i_{1},i_{2}) =  \delta_{I_{1}}(i_{3},i_{4}) \hdots = \delta_{I_{m}}(i_{k-1},i_{k}) &  i_{1}, i_{2}, \hdots , i_{k} \> \text{are pairwise distinct} \\    
  0 & \text{otherwise}    
\end{cases} \]

where $I_{1}, \hdots I_{m}$ are compatible sets of size $k-2$ satisfying compatibility conditions described in Definition \ref{def:gen_metric_tr_arr}. With the proof of Theorem \ref{prop:gen_tree_dressian} we see that for a compatible set $I \in$ $\binom{[n]}{k}$, the minimum of 

\[ \text{min} \{ \pi_{Ipq} + \pi_{Irs} , \pi_{Ipr} + \pi_{Isq} , \pi_{Ips} + \pi_{Iqr} \} \]

is attained twice, where $p,q,r,s \in [n] \setminus I$, as this minimum corresponds to the weight induced by a metric tree on $n-k+2$ leaves. Hence, we can conclude that the restriction of the map $\pi$ to increasing tuple $i_{1} < i_{2} < \hdots < i_{k}$ is a finite tropical Pl\"ucker vector and is an element of $\Dr(k,n)$.


There is a natural fan structure on the space of weights induced by generalized metric tree arrangements. Two weights $w_T$ and $w_{T'}$ lie in the interior of the same cone if and only if $T$ and $T'$ have the same underlying set of labeled trees. The proof above naturally extends and shows that the Pl\"ucker fan structure of the Dressian coincides with the fan structure on the space of weights induced by generalized metric tree arrangements. Also, as a consequence of Pl\"ucker fan structure and secondary fan structure coinciding for the Dressian \cite{olarte2019local}, the fan structure on the space of weights induced by generalized metric tree arrangements coincides with the secondary fan structure of Dressian. 

The notion of a \textit{cherry} of an arrangement naturally extends to the generalized setting. A \textit{cherry} of an arrangement $T \in \TT(k,n)$ is a $k$-subset $I \in [n]$ such that if $I = J \cup \{i,j\}$ is any partition then $i,j$ is a cherry in tree $T_J$.
\bigskip

With Theorem \ref{prop:gen_tree_dressian} we see that a generalized metric tree arrangement corresponds to a regular matroidal subdivision of $\Delta(k,n)$. We see this when we consider the collection of weights $\{w_T \colon T \in \TT(k,n) \}$ as a height function on the Pl\"ucker coordinates. Using the fact that the Pl\"ucker fan structure and secondary fan structure coincide for the Dressian, we obtain the following result,

\begin{corollary}\label{lem:sub_to_gen_trees}
Each regular matroidal subdivision $\Sigma$ of $\Delta(k,n)$ supports a generalized metric tree arrangement.  
\end{corollary}


\begin{proof}
Let us consider the weight function for the regular matroidal subdivision $\Sigma$ to be $w_{\Sigma}$ and $w_{\Sigma} \in \Dr(k,n)$. Since the Pl\"ucker fan structure and secondary fan structure coincide for the Dressian, therefore by Theorem \ref{prop:gen_tree_dressian}, we can conclude that there exists a generalized metric tree arrangement $T(k,n)$ such that it induces $w_{\Sigma}$.
\end{proof}

%\begin{corollary}
%A generalized metric tree arrangement of $n \choose k-2$ trees is equivalent to a metric tree arrangement of $\Myperm[n]{k-2}$ trees.  
%\end{corollary}

We know that in \cite{herrmann2009draw} Lemma 4.2 and Theorem 4.5 help establish a bijection between equivalence classes of \emph{equivalent} metric tree arrangements of $n$ trees and regular matroidal subdivisions of $\Delta(3,n)$. We now try to establish this equivalence for generalized metric tree arrangements. We already see a mapping from one side, i.e., by Theorem \ref{prop:gen_tree_dressian} we see that each generalized metric tree arrangement corresponds to a regular matroidal subdivision of $\Delta(k,n)$. We now have to show a mapping in the reverse direction. 

We refer to a generalized metric tree arrangement as an \emph{abstract} generalized metric tree arrangement if we consider the tree arrangement without the metric. Two generalized metric tree arrangements are \emph{equivalent} if they induce the same abstract tree arrangement. With these results, we can now state the following result which in essence is a generalization of Theorem 4.4 in \cite{herrmann2009draw}.

\begin{corollary}
The equivalence classes of arrangements of $\binom{n}{k-2}$ generalized metric tree are in bijection with the regular matroid subdivisions of $\Delta(k,n)$.
\end{corollary}

\begin{proof}
With Theorem \ref{prop:gen_tree_dressian}, we see that a generalized metric tree arrangement defines a point in the Dressian $\Dr(k,n)$. We observe that equivalent generalized metric tree arrangements define points in the $\Dr(k,n)$ which lie in the same cone in the Pl\"ucker fan. Considering this weight vector as a lifting function provides us the corresponding regular matroid subdivision. This provides us with a map from equivalence classes of generalized metric tree arrangements to regular matroidal subdivisions. For the map from the other side we consider a regular matroidal subdivision of $\Dr(k,n)$ and by Corollary \ref{lem:sub_to_gen_trees} we see that this subdivision corresponds to a generalized metric tree arrangement. By \cite{olarte2019local}, we know that the secondary fan structure and the Pl\"ucker fan structure coincides for $\Dr(k,n)$. Hence, this gives us the required bijection.
\end{proof}


\begin{example}\label{eg:Ex_Dr(4,6)}
    

We first show an example of a generalized metric tree arrangement for a point in $\Dr(4,6)$. In this case we realize that the generalized metric tree arrangement is an arrangement of $\binom{6}{4} = 15$ metric trees $T_{I}$, where $I \subseteq [6], \ |I| = 4-2 = 2$ and these trees are illustrated in Figure \ref{fig:Dr(4,6)_Ex}. Additionally, the metric on each tree $T_
{I}$ satisfies the compatibility condition described in Equation \ref{eq:comptability} which provides us the following relations 

\[
\delta_{12}(3,4) = \delta_{13}(2,4) = \delta_{14}(2,3) = \delta_{23}(1,4) = \delta_{24}(1,3) = \delta_{34}(1,2)
\]
\[
\delta_{12}(3,5) = \delta_{13}(2,5) = \delta_{15}(2,3) = \delta_{23}(1,5) = \delta_{25}(1,3) = \delta_{35}(1,2)
\]
\[
\delta_{12}(3,6) = \delta_{13}(2,6) = \delta_{16}(2,3) = \delta_{23}(1,6) = \delta_{26}(1,3) = \delta_{36}(1,2)
\]
\[
\delta_{12}(4,5) = \delta_{14}(2,5) = \delta_{15}(2,4) = \delta_{24}(1,5) = \delta_{25}(1,4) = \delta_{45}(1,2)
\]
\[
\delta_{12}(4,6) = \delta_{14}(2,6) = \delta_{16}(2,4) = \delta_{24}(1,6) = \delta_{26}(1,4) = \delta_{46}(1,2)
\]
\[
\delta_{12}(5,6) = \delta_{15}(2,6) = \delta_{16}(2,5) = \delta_{25}(1,6) = \delta_{26}(1,5) = \delta_{56}(1,2)
\]
\[
\delta_{13}(4,5) = \delta_{14}(3,5) = \delta_{15}(3,4) = \delta_{34}(1,5) = \delta_{35}(1,4) = \delta_{45}(1,3)
\]
\[
\delta_{13}(4,6) = \delta_{14}(3,6) = \delta_{16}(3,4) = \delta_{34}(1,6) = \delta_{36}(1,4) = \delta_{46}(1,3)
\]
\[
\delta_{13}(5,6) = \delta_{15}(3,6) = \delta_{16}(3,5) = \delta_{35}(1,6) = \delta_{36}(1,5) = \delta_{56}(1,3)
\]
\[
\delta_{14}(5,6) = \delta_{15}(4,6) = \delta_{16}(4,5) = \delta_{45}(1,6) = \delta_{46}(1,5) = \delta_{56}(1,4)
\]
\[
\delta_{23}(4,5) = \delta_{24}(3,5) = \delta_{25}(3,4) = \delta_{34}(2,5) = \delta_{35}(2,4) = \delta_{45}(2,3)
\]
\[
\delta_{23}(4,6) = \delta_{24}(3,6) = \delta_{26}(3,4) = \delta_{34}(2,6) = \delta_{36}(2,4) = \delta_{46}(2,3)
\]
\[
\delta_{23}(5,6) = \delta_{25}(3,6) = \delta_{26}(3,5) = \delta_{35}(2,6) = \delta_{36}(2,5) = \delta_{56}(2,3)
\]
\[
\delta_{24}(5,6) = \delta_{25}(4,6) = \delta_{26}(4,5) = \delta_{45}(2,6) = \delta_{46}(2,5) = \delta_{56}(2,4)
\]
\[
\delta_{34}(5,6) = \delta_{35}(4,6) = \delta_{36}(4,5) = \delta_{45}(3,6) = \delta_{46}(3,5) = \delta_{56}(3,4)
\]

\begin{figure}
    \centering
    \includegraphics[scale=0.7]{Dr_4,6_EX.drawio.pdf}
    \caption{A generalized metric tree arrangement which corresponds to a point in $\Dr(4,6)$.}
    \label{fig:Dr(4,6)_Ex}
\end{figure}

A weight vector which satisfies all the \emph{four-point conditions}\cite[Lemma 4.3.6]{maclagan2021introduction} on each metric tree is the following, 

\[
w = \{2,1,1,1,1,1,0,0,0,0,0,0,0,0,0 \}
\]


which gives the following regular matroidal subdivision $\Sigma_{1}$ of $\Delta(4,6)$

\[ M_{1} = \{1235,1236,1245,1246,1256,1345,1346,1356,1456,2345,2346,2356,2456\}
\]
\[ M_{2} = \{1345,1346,1356,1456,2345,2346,2356,2456,3456\} \]
\[ M_{3} = \{1234,1235,1236,1245,1246,1345,1346,2345,2346\} \]

Under the canonical isomorphism between $\Delta(k,n) \cong \Delta(n-k,n)$, $\Sigma_{1}$ corresponds to the regular matroidal subdivision $\Sigma_{2}$ of $\Delta(2,6)$ given by the weight vector 

\[ w' = \{0,0,0,0,0,0,0,0,0,1,1,1,1,1,2\} \]

and $\Sigma_{2}$ is also given by a metric tree which indexes the corresponding cone in $\Dr(2,6)$ which is shown in Figure \ref{fig:Dr(2,6)_Ex}.

\begin{figure}
    \centering
    \includegraphics[scale=0.5]{Dr_2,6_examp.drawio.pdf}
    \caption{The metric tree that corresponds to the subdivision $\Sigma_{2}$ of $\Delta(2,6)$ in Example \ref{eg:Ex_Dr(4,6)}.}
    \label{fig:Dr(2,6)_Ex}
\end{figure}

\begin{remark}
This dichotomy between generalized metric tree arrangement of $\Dr(4,6)$ and a metric tree of $\Dr(2,6)$ can also be seen in between metric tree arrangements and metric trees as well, for example in the case of $\Dr(3,5)$ and $\Dr(2,5)$, where the metric tree arrangement of $5$ trees on four leaves corresponds to a metric tree on 5 leaves.    
\end{remark}
\end{example}

It is important to notice that for $\Dr(k,n), k \geq 4$ we encounter both metric tree arrangements and generalized metric tree arrangements. However, we see that the metric arrangement obtained in Lemma \ref{lem:mat_sub_k,n} consists of $(k-2)!$ $\cdot \binom{n}{k-2}$ whereas the generalized metric tree arrangement consists of $\binom{n}{k-2}$ trees.


\begin{remark}
For $k=3$, we see that $(k-2)!$ $\cdot \binom{n}{k-2}$ = $\binom{n}{k-2}$, therefore the associated metric tree arrangement and generalized metric tree arrangement coincide in this specific case.
\end{remark}

Our study leads us to the following questions which we aim to explore in subsequent work,

\begin{itemize}
    \item Compute the Dressian $\Dr(4,8)$, which could provide  computational examples to complement our results. Although at the moment the state-of-the-art in terms of such computations is $\Dr(3,8)$ \cite{herrmann2014dressians},\cite{bendle2020parallel}.
    \item We are able to provide the notion of metric tree arrangements for all values of $\Dr(k,n)$. This leads us to ask that whether we can find an equivalent definition of a tree arrangement for the case of local Dressian $\Dr(M)$, where $M$ is a matroid. In recent work \cite{BrandtSpeyer+2022+409+420} there have been advancements in techniques to compute $\Dr(M)$ for certain matroids which could be beneficial for answering this question.
\end{itemize}

We condense some of the findings in this section with this statement,

\begin{corollary}\label{cor:final}
Each maximal cone of the Dressian $\Dr(k,n), n \geq 4, k \geq 2$ supports an arrangement of $\binom{n}{k-2}$ generalized metric trees.
\end{corollary}

\paragraph{\bf Acknowledgment} I am grateful to my advisor Prof Isabel Van Driessche and also to Prof Tom De Medts for their support during the preparation of this article. I would also like to thank Prof Michael Joswig for going through earlier drafts of this work. I thank Dr. Ahmed Umer Ashraf for the discussions which helped in coming up with the examples in the text. A.~K.~Tewari has been supported by the UGent BOF grant BOF/STA/201909/038 and Ugent BOF grant BOFFF12015001101.



\bibliographystyle{abbrv} 
\bibliography{Dressian.bib}

\appendix    


\section{Appendix} \label{app:appendix}
\subsection{Metric tree arrangements for Dr(3,6)}

%A sequence $T = (T_1, \dots, T_n)$ of metric trees is called an \textit{arrangement of metric trees for $\Dr(3,n)$} \cite{herrmann2009draw} if:
%\begin{itemize}
    %\item For each $i \in [n]$ the metric tree $T_i$ has leaves indexed by $[n] \backslash \{i\}$,
    
    %\item For each subset $\{i,j,k\} \in \binom{[n]}{3}$ we have $d_i(j,k) = d_j(i,k) = d_k(i,j)$ where $d_i(j,k)$ is the distance between leaves $j$ and $k$ in tree $T_i$.
%\end{itemize}
%If $T$ is an arrangement to metric trees then the underlying sequence of labeled trees is called a sequence of compatible trees. Each maximal cell of the Dressian $\Dr(3,n)$ is indexed by a sequence of compatible trees. 

For Dressian $\Dr(3,6)$ there are, up to symmetry, $7$ different maximal cones. We now list the seven cones along with the descriptions of the metric tree arrangements and the associated regular matroidal subdivision. 

For cone number $0$ from \href{https://users-math.au.dk/jensen/Research/G3_7/grassmann3_6.html}{\color{blue}Jensen's website\color{black}}, the tree arrangement corresponding to this cone is
\begin{align*}
&T_1: C(25,4,36), \quad 
T_2: C(15,3,46), \quad
T_3: C(16,2,45), \\
&T_4: C(26,1,35), \quad
T_5: C(12,6,34), \quad
T_6: C(13,5,24)
\end{align*}
where $C(ab,c,de)$ is the caterpillar graph with cherries $ab$ and $cd$. For each tree $T_i = (a_i b_i, c_i, d_i e_i)$ we write: $s_i$ for the internal vertex of $T_i$ adjacent to $a_i$ and $b_i$; $t_i$ for the internal vertex of $T_i$ adjacent to $c_i$; and $u_i$ for the internal vertex of $T_i$ adjacent to $d_i$ and $e_i$.

The matroid subdivision induced by this cone is given by $5$ cells whose corresponding matroids are $M_A, M_B, M_C, M_D$, and $M_E$: see Figure~\ref{fig: gr36 cone 0 tree arrangement}.


\begin{figure}[H]
    \centering
    \includegraphics[scale=0.8]{tree_arrangement_gr36_cone_0.drawio.pdf}
    \caption{The tree arrangement for Cone$_{0}$ together with the maximal cells of the corresponding matroid subdivision.}
    \label{fig: gr36 cone 0 tree arrangement}
\end{figure}



\begin{definition}
Given a compatible tree arrangement $T = (T_1, \dots T_n)$ a \textit{cherry of $T$} is a subset $\{i,j,k\} \in \binom{[n]}{3}$ such that: $ij$ is a cherry in $T_k$; $ik$ is a cherry in $T_j$; and $jk$ is a cherry in $T_i$
\end{definition}

\begin{table}[]
    \centering
    \begin{tabular}{|p{2.5cm}|p{5cm}|}
\hline
Cone Index & List of Cherries \\
\hline
$C_{0}$ & $\{125\}, \{136\}, \{246\}, \{345\} $ \\
$C_{1}$ & $\{246\}, \{345\}, \{125\} $ \\
$C_{2}$ & $\{125\}, \{246\} $ \\
$C_{3}$ & $\{125\}, \{246\} $ \\
$C_{4}$ & $\{135\}, \{246\} $ \\
$C_{5}$ & $\{246\} $ \\
$C_{6}$ & $\{\} $ \\

\hline
    \end{tabular}
    \caption{List of cones from \href{https://users-math.au.dk/jensen/Research/G3_7/grassmann3_6.html}{\color{blue}Jensen's website\color{black}} and their associated cherries}
    \label{tab:cherry_Dr(3,6)}
\end{table}

For Cone $6$ from \href{https://users-math.au.dk/jensen/Research/G3_7/grassmann3_6.html}{\color{blue}Jensen's website\color{black}}, the tree arrangement is as follows,

\begin{align*}
&T_1: C(23,5,46), \quad 
T_2: C(15,3,46), \quad
T_3: C(15,2,46), \\
&T_4: C(15,6,23), \quad
T_5: C(23,1,46), \quad
T_6: C(15,4,23)
\end{align*}

\begin{figure}[H]
    \centering
    \includegraphics[scale=0.8]{tree_arrange_Dr3,6_cone_6.pdf}
    \caption{The tree arrangement for Cone$_{6}$ together with the maximal cells of the corresponding matroid subdivision.}
    \label{fig: gr36 cone 6 tree arrangement}
\end{figure}

and the tree arrangement along with the matroidal subdivision is shown in Figure \ref{fig: gr36 cone 6 tree arrangement}.

For Cone $1$ from \href{https://users-math.au.dk/jensen/Research/G3_7/grassmann3_6.html}{\color{blue}Jensen's website\color{black}}, the tree arrangement is as follows,

\begin{align*}
&T_1: C(25,6,34), \quad 
T_2: C(15,3,46), \quad
T_3: C(26,1,45), \\
&T_4: C(26,1,35), \quad
T_5: C(12,6,34), \quad
T_6: C(15,3,24)
\end{align*}


\begin{figure}[H]
    \centering
    \includegraphics[scale=0.8]{tree_arrange_Dr3,6_cone_1.pdf}
    \caption{The tree arrangement for Cone$_{1}$ together with the maximal cells of the corresponding matroid subdivision.}
    \label{fig: gr36 cone 1 tree arrangement}
\end{figure}

and the tree arrangement along with the matroidal subdivision is shown in Figure \ref{fig: gr36 cone 1 tree arrangement}.

For Cone $2$ from \href{https://users-math.au.dk/jensen/Research/G3_7/grassmann3_6.html}{\color{blue}Jensen's website\color{black}}, the tree arrangement is as follows,

\begin{align*}
&T_1: C(25,6,34), \quad 
T_2: C(15,3,46), \quad
T_3: C(15,4,26), \\
&T_4: C(15,3,26), \quad
T_5: C(12,6,34), \quad
T_6: C(15,3,24)
\end{align*}


\begin{figure}[H]
    \centering
    \includegraphics[scale=0.8]{tree_arrange_Dr3,6_cone_2.pdf}
    \caption{The tree arrangement for Cone$_{2}$ together with the maximal cells of the corresponding matroid subdivision.}
    \label{fig: gr36 cone 2 tree arrangement}
\end{figure}

and the tree arrangement along with the matroidal subdivision is shown in Figure \ref{fig: gr36 cone 2 tree arrangement}.

For Cone $3$ from \href{https://users-math.au.dk/jensen/Research/G3_7/grassmann3_6.html}{\color{blue}Jensen's website\color{black}}, the tree arrangement is as follows,

\begin{align*}
&T_1: C(25,3,46), \quad 
T_2: C(15,3,46), \quad
T_3: C(15,2,46), \\
&T_4: C(15,3,26), \quad
T_5: C(12,3,46), \quad
T_6: C(15,3,24)
\end{align*}


\begin{figure}[H]
    \centering
    \includegraphics[scale=0.8]{tree_arrange_Dr3,6_cone_3.pdf}
    \caption{The tree arrangement for Cone$_{3}$ together with the maximal cells of the corresponding matroid subdivision.}
    \label{fig: gr36 cone 3 tree arrangement}
\end{figure}

and the tree arrangement along with the matroidal subdivision is shown in Figure \ref{fig: gr36 cone 3 tree arrangement}.

For Cone $4$ from \href{https://users-math.au.dk/jensen/Research/G3_7/grassmann3_6.html}{\color{blue}Jensen's website\color{black}}, the tree arrangement is as follows,

\begin{align*}
&T_1: C(35,2,46), \quad 
T_2: C(15,3,46), \quad
T_3: C(15,2,46), \\
&T_4: C(15,3,26), \quad
T_5: C(13,2,46), \quad
T_6: C(15,3,24)
\end{align*}


\begin{figure}[H]
    \centering
    \includegraphics[scale=0.8]{tree_arrange_Dr3,6_cone_4.pdf}
    \caption{The tree arrangement for Cone$_{4}$ together with the maximal cells of the corresponding matroid subdivision.}
    \label{fig: gr36 cone 4 tree arrangement}
\end{figure}

and the tree arrangement along with the matroidal subdivision is shown in Figure \ref{fig: gr36 cone 4 tree arrangement}.

For Cone $5$ from \href{https://users-math.au.dk/jensen/Research/G3_7/grassmann3_6.html}{\color{blue}Jensen's website\color{black}}, the tree arrangement is as follows,

\begin{align*}
&T_1: C(23,5,46), \quad 
T_2: C(15,3,46), \quad
T_3: C(15,2,46), \\
&T_4: C(15,3,26), \quad
T_5: C(23,1,46), \quad
T_6: C(15,3,24)
\end{align*}


\begin{figure}[H]
    \centering
    \includegraphics[scale=0.8]{tree_arrange_Dr3,6_cone_5.pdf}
    \caption{The tree arrangement for Cone$_{5}$ together with the maximal cells of the corresponding matroid subdivision.}
    \label{fig: gr36 cone 5 tree arrangement}
\end{figure}

and the tree arrangement along with the matroidal subdivision is shown in Figure \ref{fig: gr36 cone 5 tree arrangement}.


\subsection{Adjacent cone tree arrangements in Dr(3,6)}


%\subsubsection{Adjacent cones for Cone$_{0}$}

In this section, we try to look at what happens if we consider tree arrangements of two adjacent maximal cones in $\Dr(3,6)$. For this we do a computational study using Macaulay2 \cite{M2} and consider the Cone$_{0} = \{ 42, 52, 53, 54 \}$ (missing Pl\"ucker coordinate = $\{123\}$) where the indices represent the rays indexing an interior point in the cone. We compute all the cones adjacent to Cone$_{0}$ in $\Dr(3,6)$ and compute their respective tree arrangements. Cone$_{0}$ lies in the equivalence class of Cone$_{5}$ in the list on \href{https://users-math.au.dk/jensen/Research/G3_7/grassmann3_6.html}{\color{blue}Jensen's website\color{black}}. There are eight adjacent cones to Cone$_{0}$ and Table \ref{tab:adjacent_cones_Dr(3,6)_0} lists them with the rays indexing each of them and also the missing Pl\"lucker coordinates. The missing Pl\"lucker coordinates for each cone are exactly the \emph{cherries} of the tree arrangement.   


\begin{table}[H]
    \centering
    \begin{tabular}{|p{2.5cm}|p{3cm}| p{3cm}|}
\hline
Cone Index & List of Rays & Cherries \\
\hline
$C_{9}$ & $\{15, 42, 53, 54\} $ & $\{ 123,256 \}$ \\
$C_{10}$ & $\{12, 42, 53, 54\} $ &  $\{ 123, 156\}$ \\
$C_{11}$ & $\{14, 42, 52, 54\} $ & $\{ 123, 346\}$ \\
$C_{12}$ & $\{13, 42, 52, 54\} $ & $\{ 123, 345\}$ \\
$C_{81}$ & $\{52, 53, 54, 64\} $ & $\{ 123, 456\}$ \\
$C_{90}$ & $\{20, 52, 53, 54\} $ & $\{ 123, 356\}$ \\
$C_{105}$ & $\{25, 42, 52, 53, 63\} $ & $\{\}$ \\
$C_{114}$ & $\{42, 43, 52, 53\} $ & $\{ 124\}$ \\
\hline
    \end{tabular}
    \caption{List of adjacent cones for Cone$_{0}$}
    \label{tab:adjacent_cones_Dr(3,6)_0}
\end{table}

The tree arrangement for Cone$_{0}$ is shown in Figure \ref{fig: Dr36 cone 0 adjacent cone tree arrangement}. It is obtained by first computing the contraction of all maximal cells of the matroidal subdivision $\Sigma$ for the cone with respect to each index $\{i\}$ in $[n]$. From the structure of matroidal subdivision on $\Delta(2,5)$, we know that the cell indexed by an internal vertex attached to a cherry of a tree $T_{i}$ has 7 elements, while the middle internal vertex has 8 elements. Hence, for each contraction on $\Sigma$, we obtain three unique cells $M_{a}/\{i\}, M_{b}/\{i\}$ and $M_{c}/\{i\}$, out of which two have 7 elements and one has 8 elements. If the 7-element cell indexes a vertex attached to a cherry $pq$ then all elements in the cell contain either $p$ or $q$. With this, we obtain labels for the cherry of each tree $T_{i}$. The label for the middle leaf is the one index left from $[6]$ after labeling the cherries of the tree and the tree index. 

\begin{figure}[H]
    \centering
    \includegraphics[scale=0.8]{Adjacent_cones_comp_cone_0.pdf}
    \caption{The tree arrangement for Cone$_{0}$ together with the maximal cells of the corresponding matroid subdivision.}
    \label{fig: Dr36 cone 0 adjacent cone tree arrangement}
\end{figure}

%We can also associate a tropical line arrangement to an arrangement of trees, which we discuss for the case of tree arrangement for Cone$_{0}$. We choose $ijk$ for the tree arrangement, such that $i$ and $j$ belong to the two cherries in tree $T_{k}$. The six trees in the tree arrangement define six lines , three of them are lines at infinity, while the other three define a mixed subdivision of $3 \cdot \Delta_{2}$. The lines at infinity are indexed by $ijk$ choosen earlier and they are indexed by the boundary edges of $3 \cdot \Delta_{2}$. We show  in Figure \ref{fig: Dr36 cone 0 trop line arrange} the tropical line arrangement which corresponds to the tree arrangement for Cone$_0$ and the dual mixed subdivision of $3 \cdot \Delta_{2}$. 

%\begin{figure}[H]
    %\centering
    %\includegraphics[scale=0.4]{tropical_line_arrangement_Cone_0.drawio.pdf}
    %\caption{The tropical line arrangement corresponding to the tree arrangement for $Cone_9$ %together with the mixed subdivision of $3 \cdot \Delta_{2}$.}
    %\label{fig: Dr36 cone 0 trop line arrange}
%\end{figure}

For the adjacent Cone$_{9}$, we have the tree arrangement shown in Figure \ref{fig: Dr36 cone 9 adjacent cone tree arrangement}

\begin{figure}[H]
    \centering
    \includegraphics[scale=0.7]{Adjacent_cones_comp_cone_0_9.pdf}
    \caption{The tree arrangement for Cone$_9$ together with the maximal cells of the corresponding matroid subdivision. The shaded edges represent the labels that get changed with respect to the labeling in Cone$_{0}$ }
    \label{fig: Dr36 cone 9 adjacent cone tree arrangement}
\end{figure}

For the adjacent Cone$_{10}$, we have the tree arrangement shown in Figure \ref{fig: Dr36 cone 10 adjacent cone tree arrangement}

\begin{figure}[H]
    \centering
    \includegraphics[scale=0.8]{Adjacent_cones_comp_cone_0_10.pdf}
    \caption{The tree arrangement for Cone$_{10}$ together with the maximal cells of the corresponding matroid subdivision. The shaded edges represent the labels that get changed with respect to the labeling in Cone$_{0}$ }
    \label{fig: Dr36 cone 10 adjacent cone tree arrangement}
\end{figure}

For the adjacent Cone$_{11}$, we have the tree arrangement shown in Figure \ref{fig: Dr36 cone 11 adjacent cone tree arrangement}

\begin{figure}[H]
    \centering
    \includegraphics[scale=0.8]{Adjacent_cones_comp_cone_0_11.pdf}
    \caption{The tree arrangement for Cone$_{11}$ together with the maximal cells of the corresponding matroid subdivision. The shaded edges represent the labels that get changed with respect to the labeling in Cone$_{0}$ }
    \label{fig: Dr36 cone 11 adjacent cone tree arrangement}
\end{figure}

For the adjacent Cone$_{12}$, we have the tree arrangement shown in Figure \ref{fig: Dr36 cone 12 adjacent cone tree arrangement}

\begin{figure}[H]
    \centering
    \includegraphics[scale=0.8]{Adjacent_cones_comp_cone_0_12.pdf}
    \caption{The tree arrangement for Cone$_{12}$ together with the maximal cells of the corresponding matroid subdivision. The shaded edges represent the labels that get changed with respect to the labeling in Cone$_{0}$ }
    \label{fig: Dr36 cone 12 adjacent cone tree arrangement}
\end{figure}

For the adjacent Cone$_{81}$, we have the tree arrangement shown in Figure \ref{fig: Dr36 cone 81 adjacent cone tree arrangement}

\begin{figure}[H]
    \centering
    \includegraphics[scale=0.8]{Adjacent_cones_comp_cone_0_81.pdf}
    \caption{The tree arrangement for Cone$_{81}$ together with the maximal cells of the corresponding matroid subdivision. The shaded edges represent the labels that get changed with respect to the labeling in Cone$_{0}$ }
    \label{fig: Dr36 cone 81 adjacent cone tree arrangement}
\end{figure}

For the adjacent Cone$_{90}$, we have the tree arrangement shown in Figure \ref{fig: Dr36 cone 90 adjacent cone tree arrangement}

\begin{figure}[H]
    \centering
    \includegraphics[scale=0.8]{Adjacent_cones_comp_cone_0_90.pdf}
    \caption{The tree arrangement for Cone$_{90}$ together with the maximal cells of the corresponding matroid subdivision. The shaded edges represent the labels that get changed with respect to the labeling in Cone$_{0}$ }
    \label{fig: Dr36 cone 90 adjacent cone tree arrangement}
\end{figure}

For the adjacent Cone$_{105}$, we have the tree arrangement shown in Figure \ref{fig: Dr36 cone 105 adjacent cone tree arrangement}

\begin{figure}[H]
    \centering
    \includegraphics[scale=0.8]{Adjacent_cones_comp_cone_0_105.pdf}
    \caption{The tree arrangement for Cone$_{105}$ together with the maximal cells of the corresponding matroid subdivision. The shaded edges represent the labels that get changed with respect to the labeling in Cone$_{0}$ }
    \label{fig: Dr36 cone 105 adjacent cone tree arrangement}
\end{figure}


For the adjacent Cone$_{114}$, we have the tree arrangement shown in Figure \ref{fig: Dr36 cone 114 adjacent cone tree arrangement}

\begin{figure}[H]
    \centering
    \includegraphics[scale=0.8]{Adjacent_cones_comp_cone_0_114.pdf}
    \caption{The tree arrangement for Cone$_{114}$ together with the maximal cells of the corresponding matroid subdivision. The shaded edges represent the labels that get changed with respect to the labeling in Cone$_{0}$ }
    \label{fig: Dr36 cone 114 adjacent cone tree arrangement}
\end{figure}

%We now try to look at the mixed subdivisions of $3 \cdot \Delta_{2}$ for all adjacent cones of Cone$_{0}$.

%\begin{figure}
    %\centering
    %\includegraphics[scale=0.35]{all_subdivision_adjacent.drawio.pdf}
    %\caption{The mixed subdivision of $3 \cdot \Delta_{2}$ for all adjacent cones of Cone$_{0}$. }
    %\label{fig: adjacent cones Dr36 mixed subdiv}
%\end{figure}

Hence, we see that for Cone$_{0}$, adjacent cones have tree arrangements wherein all trees are identical except two trees which are related by a Whitehead move. We conduct this computational study for cones of every equivalence class out of the seven distinct classes of cones for $\Dr(3,6)$. We observe that in all of them, for adjacent cones, all trees are identical except at most three trees which are related via a Whitehead move. The computational files for the other six equivalence classes can be found at 

\begin{center}
    
\url{https://github.com/Ayush-Tewari13/Cone_Adjacency_Dr-3-6-}

\end{center}




\end{document}
